\documentclass[12pt, a4paper]{article}

% Paquetes requeridos para el circuito
\usepackage[utf8]{inputenc}
\usepackage[T1]{fontenc}
\usepackage{circuitikz} 

\begin{document}

\begin{center}
    \begin{circuitikz}[american, scale=1.0, transform shape]
        
        % 1. Raíl inferior común (nodo 'b')
        \draw (0,0) -- (8,0) node[circ] {} node[right] {b};
        
        % 2. Fuente independiente de tensión Vs entre 'b' y 'c'
        \draw (0,0) node[circ] {} to[V, l=$V_s$] (0,3) node[circ] {} node[above] {c};
        
        % 3. Resistencia Rg horizontal entre 'c' y 'd'
        \draw (0,3) to[R, l=$R_g$] (3,3) node[circ] {} node[above] {d};
        
        % 4. Resistencia Ri entre 'd' y 'b', con la tensión de control V1
        \draw (3,3) to[R, l=$R_i$, v=$V_1$] (3,0) node[circ] {};
        
        % 5. Etapa de salida: Fuente dependiente de tensión K*V1 entre 'b' y 'e'
        \draw (5.5,0) node[circ] {} to[cV, l=$K V_1$, invert] (5.5,3) node[circ] {} node[above] {e};
        
        % 6. Conexión superior de la etapa de salida (de 'e' a 'a') y RL con tensión Vo
        \draw (5.5,3) to[R, l=$R_o$] (8,3) node[circ] {} node[above] {a};
        \draw (8,3) to[R, l=$R_L$, v=$V_o$] (8,0) node[circ] {};
        
    \end{circuitikz}
\end{center}

\end{document}